\clearpage\phantomsection\label{couverture-programme}\pdfbookmark[0]{Correspondance avec le programme}{bm-couverture-programme}
\gdef\otChapNum{0}\pagestyle{otFront}\otPageTitre{Programme et chapitres}{Troisième · Programme 2026 applicable à la rentrée 2028}
Chaque rubrique de Troisième de l’annexe 2 de l’arrêté du 18 février 2026 est
reliée aux chapitres qui la traitent. Les pages suivantes associent ensuite quatre
repères par chapitre à des exercices précis ; les séries citées ne sont pas exhaustives.\par\medskip
{\fontsize{9.5}{12.5}\selectfont
\begin{tabularx}{\linewidth}{@{}>{\raggedright\arraybackslash}p{32mm}Yr@{}}\toprule
Domaine & Rubrique de Troisième & Chapitres\\\midrule
Nombres et calculs & Nombres rationnels : forme irréductible, problèmes ; notation des ensembles de nombres (prolongement) & 1\\[1.6mm]
 & Puissances : exposants négatifs, produits et quotients, notation scientifique & 2, 7\\[1.6mm]
 & Racine carrée : équations $x^2=a$, résolution par le calcul et graphique, problèmes & 2, 5, 8\\[1.6mm]
 & Multiples et diviseurs : décomposition, critères de divisibilité, fractions (automatismes) & 3, 1\\[1.6mm]
 & Calcul littéral : produits et rapports, double distributivité, identités remarquables, produit nul, inéquation $ax\ge b$, analyse-synthèse & 4, 5\\[1.6mm]
Espace et géométrie & Repérage sur une droite graduée et dans le plan (automatismes) & 5, 6, 8, 9\\[1.6mm]
 & Représentation de l’espace : solides et volumes, boule, sphère, grands cercles, sections planes & 16\\[1.6mm]
 & Triangles : Thalès, réciproque et contraposée ; cosinus, sinus, tangente ; Pythagore, cercle circonscrit et droite des milieux (automatismes) & 12, 13, 14\\[1.6mm]
 & Translations et vecteurs : vecteurs égaux, nul, opposé, somme, relation de Chasles & 15\\[1.6mm]
Données et probabilités & Statistiques : effectifs cumulés, médiane, quartiles, boîtes à moustaches, tableur & 10\\[1.6mm]
 & Probabilités : relation $P(A\cup B)+P(A\cap B)=P(A)+P(B)$, expériences indépendantes, simulation et stabilisation des fréquences & 11, 17\\[1.6mm]
Proportionnalité, fonctions & Proportionnalité : ratios, pourcentages, échelles, coefficient multiplicateur, lien avec Thalès, fonctions linéaires & 6, 7, 9\\[1.6mm]
 & Fonctions : représentations, image et antécédents, fonctions linéaires et affines, inéquations, fonction carré & 8, 9\\[1.6mm]
Pensée informatique & Variables, conditions composées, boucle conditionnelle, blocs personnalisés, écriture de programmes & 17\\[1.6mm]
Principes du programme & Raisonnement : conjecture, preuve, contre-exemple, réciproque, contraposée, absurde ; automatismes et résolution de problèmes & 18 et tous\\[1.6mm]
\bottomrule\end{tabularx}}\par\medskip
\begin{attention}Les automatismes de Troisième s’appuient sur des notions étudiées en Quatrième sans y être automatisées. L’épreuve évalue les connaissances et compétences attendues en fin de cycle 4, déclinées par le programme de Troisième. Les défis prolongent ces notions sans ajouter d’exigence à l’examen.\end{attention}
\clearpage
\otPageTitre{Repères et exercices}{Chapitres 1 à 3}\fontsize{10}{13}\selectfont
\otIntertitre{1. Nombres rationnels et fractions}\begin{tabularx}{\linewidth}{@{}Yl@{}}\toprule
Priorités, signes et vocabulaire & \hyperref[exo-1-1]{1.1}, \hyperref[exo-1-2]{1.2}, \hyperref[exo-1-10]{1.10}\\[2mm]
Opérations et fractions irréductibles & \hyperref[exo-1-3]{1.3}, \hyperref[exo-1-4]{1.4}, \hyperref[exo-1-5]{1.5}, \hyperref[exo-1-6]{1.6}, \hyperref[exo-1-8]{1.8}\\[2mm]
Comparaisons et ensembles de nombres & \hyperref[exo-1-7]{1.7}, \hyperref[exo-1-9]{1.9}, \hyperref[exo-1-15]{1.15}\\[2mm]
Fractions dans les problèmes & \hyperref[exo-1-11]{1.11}, \hyperref[exo-1-12]{1.12}, \hyperref[exo-1-13]{1.13}, \hyperref[exo-1-14]{1.14}\\[2mm]
\bottomrule\end{tabularx}\par\vspace{4mm}
\otIntertitre{2. Puissances, racines et notation scientifique}\begin{tabularx}{\linewidth}{@{}Yl@{}}\toprule
Puissances d’exposant entier & \hyperref[exo-2-1]{2.1}, \hyperref[exo-2-2]{2.2}, \hyperref[exo-2-3]{2.3}, \hyperref[exo-2-11]{2.11}\\[2mm]
Notation scientifique et préfixes & \hyperref[exo-2-4]{2.4}, \hyperref[exo-2-6]{2.6}, \hyperref[exo-2-7]{2.7}, \hyperref[exo-2-8]{2.8}, \hyperref[exo-2-12]{2.12}, \hyperref[exo-2-13]{2.13}\\[2mm]
Racine carrée et équation $x^2=a$ & \hyperref[exo-2-5]{2.5}, \hyperref[exo-2-9]{2.9}, \hyperref[exo-2-10]{2.10}, \hyperref[exo-2-15]{2.15}\\[2mm]
Problèmes de grandeurs et de racines & \hyperref[exo-2-11]{2.11}, \hyperref[exo-2-12]{2.12}, \hyperref[exo-2-13]{2.13}, \hyperref[exo-2-14]{2.14}, \hyperref[exo-2-15]{2.15}\\[2mm]
\bottomrule\end{tabularx}\par\vspace{4mm}
\otIntertitre{3. Multiples, diviseurs et nombres premiers}\begin{tabularx}{\linewidth}{@{}Yl@{}}\toprule
Divisibilité et nombres premiers & \hyperref[exo-3-1]{3.1}, \hyperref[exo-3-2]{3.2}, \hyperref[exo-3-5]{3.5}, \hyperref[exo-3-6]{3.6}, \hyperref[exo-3-7]{3.7}, \hyperref[exo-3-10]{3.10}\\[2mm]
Décomposition et fractions irréductibles & \hyperref[exo-3-3]{3.3}, \hyperref[exo-3-4]{3.4}, \hyperref[exo-3-8]{3.8}, \hyperref[exo-3-9]{3.9}\\[2mm]
Diviseurs communs et partages & \hyperref[exo-3-11]{3.11}, \hyperref[exo-3-13]{3.13}\\[2mm]
Multiples communs et périodicité & \hyperref[exo-3-12]{3.12}, \hyperref[exo-3-14]{3.14}, \hyperref[exo-3-15]{3.15}\\[2mm]
\bottomrule\end{tabularx}\par\vspace{4mm}
\clearpage
\otPageTitre{Repères et exercices}{Chapitres 4 à 6}\fontsize{10}{13}\selectfont
\otIntertitre{4. Calcul littéral et identités remarquables}\begin{tabularx}{\linewidth}{@{}Yl@{}}\toprule
Nature, opposé et substitution & \hyperref[exo-4-1]{4.1}, \hyperref[exo-4-2]{4.2}, \hyperref[exo-4-3]{4.3}\\[2mm]
Simplifier et développer & \hyperref[exo-4-4]{4.4}, \hyperref[exo-4-6]{4.6}, \hyperref[exo-4-8]{4.8}, \hyperref[exo-4-12]{4.12}\\[2mm]
Identités remarquables et factorisation & \hyperref[exo-4-5]{4.5}, \hyperref[exo-4-7]{4.7}, \hyperref[exo-4-9]{4.9}, \hyperref[exo-4-10]{4.10}\\[2mm]
Preuves par le calcul littéral & \hyperref[exo-4-11]{4.11}, \hyperref[exo-4-13]{4.13}, \hyperref[exo-4-14]{4.14}, \hyperref[exo-4-15]{4.15}\\[2mm]
\bottomrule\end{tabularx}\par\vspace{4mm}
\otIntertitre{5. Équations, inéquations et produits nuls}\begin{tabularx}{\linewidth}{@{}Yl@{}}\toprule
Équations du premier degré & \hyperref[exo-5-1]{5.1}, \hyperref[exo-5-2]{5.2}, \hyperref[exo-5-3]{5.3}, \hyperref[exo-5-4]{5.4}, \hyperref[exo-5-6]{5.6}, \hyperref[exo-5-7]{5.7}\\[2mm]
Inéquations et représentation des solutions & \hyperref[exo-5-5]{5.5}, \hyperref[exo-5-10]{5.10}, \hyperref[exo-5-13]{5.13}\\[2mm]
Produits nuls et factorisations & \hyperref[exo-5-8]{5.8}, \hyperref[exo-5-9]{5.9}, \hyperref[exo-5-15]{5.15}\\[2mm]
Mise en équation, analyse et synthèse & \hyperref[exo-5-11]{5.11}, \hyperref[exo-5-12]{5.12}, \hyperref[exo-5-13]{5.13}, \hyperref[exo-5-14]{5.14}, \hyperref[exo-5-15]{5.15}\\[2mm]
\bottomrule\end{tabularx}\par\vspace{4mm}
\otIntertitre{6. Proportionnalité, ratios et pourcentages}\begin{tabularx}{\linewidth}{@{}Yl@{}}\toprule
Proportionnalité, coefficient et alignement & \hyperref[exo-6-1]{6.1}, \hyperref[exo-6-2]{6.2}, \hyperref[exo-6-15]{6.15}\\[2mm]
Ratios, partages et échelles & \hyperref[exo-6-5]{6.5}, \hyperref[exo-6-8]{6.8}, \hyperref[exo-6-9]{6.9}, \hyperref[exo-6-10]{6.10}\\[2mm]
Évolutions et coefficient multiplicateur & \hyperref[exo-6-3]{6.3}, \hyperref[exo-6-4]{6.4}, \hyperref[exo-6-6]{6.6}, \hyperref[exo-6-7]{6.7}\\[2mm]
Comparaisons et pourcentages emboîtés & \hyperref[exo-6-11]{6.11}, \hyperref[exo-6-12]{6.12}, \hyperref[exo-6-13]{6.13}, \hyperref[exo-6-14]{6.14}\\[2mm]
\bottomrule\end{tabularx}\par\vspace{4mm}
\clearpage
\otPageTitre{Repères et exercices}{Chapitres 7 à 9}\fontsize{10}{13}\selectfont
\otIntertitre{7. Grandeurs quotients et conversions}\begin{tabularx}{\linewidth}{@{}Yl@{}}\toprule
Conversions de longueurs, aires et volumes & \hyperref[exo-7-1]{7.1}, \hyperref[exo-7-2]{7.2}, \hyperref[exo-7-13]{7.13}, \hyperref[exo-7-15]{7.15}\\[2mm]
Durées et vitesses moyennes & \hyperref[exo-7-3]{7.3}, \hyperref[exo-7-4]{7.4}, \hyperref[exo-7-6]{7.6}, \hyperref[exo-7-7]{7.7}, \hyperref[exo-7-11]{7.11}\\[2mm]
Débits et masses volumiques & \hyperref[exo-7-5]{7.5}, \hyperref[exo-7-8]{7.8}, \hyperref[exo-7-12]{7.12}, \hyperref[exo-7-14]{7.14}, \hyperref[exo-7-15]{7.15}\\[2mm]
Grandeurs quotients et contrôle des unités & \hyperref[exo-7-9]{7.9}, \hyperref[exo-7-10]{7.10}, \hyperref[exo-7-13]{7.13}, \hyperref[exo-7-14]{7.14}\\[2mm]
\bottomrule\end{tabularx}\par\vspace{4mm}
\otIntertitre{8. Fonctions, images et fonction carré}\begin{tabularx}{\linewidth}{@{}Yl@{}}\toprule
Vocabulaire, notations et calcul d’images & \hyperref[exo-8-1]{8.1}, \hyperref[exo-8-3]{8.3}, \hyperref[exo-8-5]{8.5}, \hyperref[exo-8-6]{8.6}\\[2mm]
Tableaux, courbes et antécédents & \hyperref[exo-8-2]{8.2}, \hyperref[exo-8-4]{8.4}, \hyperref[exo-8-7]{8.7}, \hyperref[exo-8-9]{8.9}, \hyperref[exo-8-14]{8.14}\\[2mm]
Fonction carré et équation $x^2=a$ & \hyperref[exo-8-8]{8.8}, \hyperref[exo-8-10]{8.10}, \hyperref[exo-8-12]{8.12}\\[2mm]
Modèles et domaines de validité & \hyperref[exo-8-11]{8.11}, \hyperref[exo-8-13]{8.13}, \hyperref[exo-8-14]{8.14}, \hyperref[exo-8-15]{8.15}\\[2mm]
\bottomrule\end{tabularx}\par\vspace{4mm}
\otIntertitre{9. Fonctions linéaires et affines}\begin{tabularx}{\linewidth}{@{}Yl@{}}\toprule
Fonctions linéaires et proportionnalité & \hyperref[exo-9-1]{9.1}, \hyperref[exo-9-2]{9.2}, \hyperref[exo-9-4]{9.4}, \hyperref[exo-9-6]{9.6}, \hyperref[exo-9-13]{9.13}\\[2mm]
Fonctions affines et tracés & \hyperref[exo-9-5]{9.5}, \hyperref[exo-9-10]{9.10}, \hyperref[exo-9-12]{9.12}, \hyperref[exo-9-14]{9.14}\\[2mm]
Coefficients lus et calculés & \hyperref[exo-9-3]{9.3}, \hyperref[exo-9-7]{9.7}, \hyperref[exo-9-8]{9.8}, \hyperref[exo-9-9]{9.9}\\[2mm]
Inéquations, comparaisons et seuils & \hyperref[exo-9-11]{9.11}, \hyperref[exo-9-13]{9.13}, \hyperref[exo-9-15]{9.15}\\[2mm]
\bottomrule\end{tabularx}\par\vspace{4mm}
\clearpage
\otPageTitre{Repères et exercices}{Chapitres 10 à 12}\fontsize{10}{13}\selectfont
\otIntertitre{10. Statistiques et boîtes à moustaches}\begin{tabularx}{\linewidth}{@{}Yl@{}}\toprule
Moyenne, médiane et étendue & \hyperref[exo-10-1]{10.1}, \hyperref[exo-10-2]{10.2}, \hyperref[exo-10-4]{10.4}, \hyperref[exo-10-10]{10.10}, \hyperref[exo-10-11]{10.11}\\[2mm]
Effectifs cumulés et quartiles & \hyperref[exo-10-3]{10.3}, \hyperref[exo-10-6]{10.6}, \hyperref[exo-10-7]{10.7}, \hyperref[exo-10-9]{10.9}\\[2mm]
Boîtes à moustaches & \hyperref[exo-10-5]{10.5}, \hyperref[exo-10-8]{10.8}, \hyperref[exo-10-12]{10.12}\\[2mm]
Comparer et interpréter des séries & \hyperref[exo-10-12]{10.12}, \hyperref[exo-10-13]{10.13}, \hyperref[exo-10-14]{10.14}, \hyperref[exo-10-15]{10.15}\\[2mm]
\bottomrule\end{tabularx}\par\vspace{4mm}
\otIntertitre{11. Probabilités, évènements et simulations}\begin{tabularx}{\linewidth}{@{}Yl@{}}\toprule
Évènements, contraire et notations & \hyperref[exo-11-1]{11.1}, \hyperref[exo-11-2]{11.2}, \hyperref[exo-11-3]{11.3}, \hyperref[exo-11-4]{11.4}, \hyperref[exo-11-13]{11.13}\\[2mm]
Réunion, intersection et « ou » exclusif & \hyperref[exo-11-6]{11.6}, \hyperref[exo-11-9]{11.9}, \hyperref[exo-11-10]{11.10}, \hyperref[exo-11-14]{11.14}\\[2mm]
Dénombrer deux épreuves indépendantes & \hyperref[exo-11-7]{11.7}, \hyperref[exo-11-8]{11.8}, \hyperref[exo-11-10]{11.10}, \hyperref[exo-11-11]{11.11}, \hyperref[exo-11-14]{11.14}\\[2mm]
Fréquences et simulations & \hyperref[exo-11-5]{11.5}, \hyperref[exo-11-12]{11.12}, \hyperref[exo-11-15]{11.15}\\[2mm]
\bottomrule\end{tabularx}\par\vspace{4mm}
\otIntertitre{12. Pythagore, cercle et droite des milieux}\begin{tabularx}{\linewidth}{@{}Yl@{}}\toprule
Pythagore : calculer une longueur & \hyperref[exo-12-1]{12.1}, \hyperref[exo-12-2]{12.2}, \hyperref[exo-12-3]{12.3}, \hyperref[exo-12-4]{12.4}, \hyperref[exo-12-11]{12.11}, \hyperref[exo-12-14]{12.14}\\[2mm]
Réciproque et contraposée & \hyperref[exo-12-6]{12.6}, \hyperref[exo-12-7]{12.7}, \hyperref[exo-12-10]{12.10}\\[2mm]
Cercle circonscrit au triangle rectangle & \hyperref[exo-12-5]{12.5}, \hyperref[exo-12-12]{12.12}, \hyperref[exo-12-13]{12.13}\\[2mm]
Théorèmes des milieux & \hyperref[exo-12-8]{12.8}, \hyperref[exo-12-9]{12.9}, \hyperref[exo-12-15]{12.15}\\[2mm]
\bottomrule\end{tabularx}\par\vspace{4mm}
\clearpage
\otPageTitre{Repères et exercices}{Chapitres 13 à 15}\fontsize{10}{13}\selectfont
\otIntertitre{13. Thalès, réciproque et contraposée}\begin{tabularx}{\linewidth}{@{}Yl@{}}\toprule
Thalès emboîté et en papillon & \hyperref[exo-13-1]{13.1}, \hyperref[exo-13-2]{13.2}, \hyperref[exo-13-3]{13.3}, \hyperref[exo-13-4]{13.4}, \hyperref[exo-13-5]{13.5}, \hyperref[exo-13-8]{13.8}, \hyperref[exo-13-9]{13.9}, \hyperref[exo-13-13]{13.13}\\[2mm]
Réciproque du théorème de Thalès & \hyperref[exo-13-6]{13.6}, \hyperref[exo-13-10]{13.10}, \hyperref[exo-13-14]{13.14}\\[2mm]
Contraposée et non-parallélisme & \hyperref[exo-13-7]{13.7}, \hyperref[exo-13-15]{13.15}\\[2mm]
Modélisation et preuve par les aires & \hyperref[exo-13-11]{13.11}, \hyperref[exo-13-12]{13.12}, \hyperref[exo-13-14]{13.14}\\[2mm]
\bottomrule\end{tabularx}\par\vspace{4mm}
\otIntertitre{14. Trigonométrie du triangle rectangle}\begin{tabularx}{\linewidth}{@{}Yl@{}}\toprule
Choisir cosinus, sinus ou tangente & \hyperref[exo-14-1]{14.1}, \hyperref[exo-14-2]{14.2}, \hyperref[exo-14-3]{14.3}, \hyperref[exo-14-4]{14.4}, \hyperref[exo-14-15]{14.15}\\[2mm]
Calculer une longueur & \hyperref[exo-14-5]{14.5}, \hyperref[exo-14-6]{14.6}, \hyperref[exo-14-7]{14.7}, \hyperref[exo-14-13]{14.13}, \hyperref[exo-14-14]{14.14}\\[2mm]
Calculer un angle aigu & \hyperref[exo-14-8]{14.8}, \hyperref[exo-14-9]{14.9}, \hyperref[exo-14-12]{14.12}\\[2mm]
Relation $\cos^2+\sin^2=1$ et mesures indirectes & \hyperref[exo-14-10]{14.10}, \hyperref[exo-14-11]{14.11}, \hyperref[exo-14-14]{14.14}\\[2mm]
\bottomrule\end{tabularx}\par\vspace{4mm}
\otIntertitre{15. Translations et vecteurs}\begin{tabularx}{\linewidth}{@{}Yl@{}}\toprule
Symétrie, demi-tour et translation & \hyperref[exo-15-1]{15.1}, \hyperref[exo-15-4]{15.4}, \hyperref[exo-15-10]{15.10}, \hyperref[exo-15-14]{15.14}\\[2mm]
Vecteurs égaux, nul et opposé & \hyperref[exo-15-2]{15.2}, \hyperref[exo-15-3]{15.3}, \hyperref[exo-15-5]{15.5}, \hyperref[exo-15-12]{15.12}\\[2mm]
Parallélogramme et égalité de vecteurs & \hyperref[exo-15-4]{15.4}, \hyperref[exo-15-6]{15.6}, \hyperref[exo-15-7]{15.7}, \hyperref[exo-15-11]{15.11}, \hyperref[exo-15-15]{15.15}\\[2mm]
Somme de vecteurs et relation de Chasles & \hyperref[exo-15-8]{15.8}, \hyperref[exo-15-9]{15.9}, \hyperref[exo-15-10]{15.10}, \hyperref[exo-15-11]{15.11}, \hyperref[exo-15-13]{15.13}\\[2mm]
\bottomrule\end{tabularx}\par\vspace{4mm}
\clearpage
\otPageTitre{Repères et exercices}{Chapitres 16 à 18}\fontsize{10}{13}\selectfont
\otIntertitre{16. Espace : sphère, boule et sections}\begin{tabularx}{\linewidth}{@{}Yl@{}}\toprule
Volumes usuels et boule & \hyperref[exo-16-1]{16.1}, \hyperref[exo-16-2]{16.2}, \hyperref[exo-16-3]{16.3}, \hyperref[exo-16-4]{16.4}, \hyperref[exo-16-5]{16.5}, \hyperref[exo-16-7]{16.7}\\[2mm]
Sphère, boule et grands cercles & \hyperref[exo-16-5]{16.5}, \hyperref[exo-16-6]{16.6}, \hyperref[exo-16-10]{16.10}\\[2mm]
Sections planes de solides & \hyperref[exo-16-8]{16.8}, \hyperref[exo-16-9]{16.9}, \hyperref[exo-16-10]{16.10}, \hyperref[exo-16-13]{16.13}, \hyperref[exo-16-14]{16.14}\\[2mm]
Assemblages et problèmes de volumes & \hyperref[exo-16-11]{16.11}, \hyperref[exo-16-12]{16.12}, \hyperref[exo-16-13]{16.13}, \hyperref[exo-16-15]{16.15}\\[2mm]
\bottomrule\end{tabularx}\par\vspace{4mm}
\otIntertitre{17. Pensée informatique et algorithmique}\begin{tabularx}{\linewidth}{@{}Yl@{}}\toprule
Variables et exécution à la main & \hyperref[exo-17-1]{17.1}, \hyperref[exo-17-2]{17.2}, \hyperref[exo-17-6]{17.6}, \hyperref[exo-17-10]{17.10}\\[2mm]
Conditions simples et composées & \hyperref[exo-17-4]{17.4}, \hyperref[exo-17-8]{17.8}, \hyperref[exo-17-9]{17.9}, \hyperref[exo-17-14]{17.14}\\[2mm]
Boucles bornées et conditionnelles & \hyperref[exo-17-3]{17.3}, \hyperref[exo-17-5]{17.5}, \hyperref[exo-17-7]{17.7}, \hyperref[exo-17-10]{17.10}, \hyperref[exo-17-11]{17.11}\\[2mm]
Blocs personnalisés et simulations & \hyperref[exo-17-12]{17.12}, \hyperref[exo-17-13]{17.13}, \hyperref[exo-17-14]{17.14}, \hyperref[exo-17-15]{17.15}\\[2mm]
\bottomrule\end{tabularx}\par\vspace{4mm}
\otIntertitre{18. Raisonner et préparer le brevet}\begin{tabularx}{\linewidth}{@{}Yl@{}}\toprule
Conjecture, preuve et contre-exemple & \hyperref[exo-18-1]{18.1}, \hyperref[exo-18-2]{18.2}, \hyperref[exo-18-6]{18.6}, \hyperref[exo-18-8]{18.8}, \hyperref[exo-18-10]{18.10}\\[2mm]
Réciproque, contraposée et absurde & \hyperref[exo-18-3]{18.3}, \hyperref[exo-18-5]{18.5}, \hyperref[exo-18-7]{18.7}\\[2mm]
Problèmes de synthèse & \hyperref[exo-18-11]{18.11}, \hyperref[exo-18-12]{18.12}, \hyperref[exo-18-13]{18.13}, \hyperref[exo-18-14]{18.14}\\[2mm]
Méthode et format du brevet 2029 & \hyperref[exo-18-4]{18.4}, \hyperref[exo-18-9]{18.9}, \hyperref[exo-18-15]{18.15}\\[2mm]
\bottomrule\end{tabularx}\par\vspace{4mm}
\clearpage
